\documentclass[11pt]{article}
\usepackage[margin=2.5cm]{geometry}
\usepackage{amsmath,amssymb}
\usepackage{xcolor}
\newcommand{\fillin}[1]{\textcolor{red}{[#1]}}

\title{Validation Plan: Harmonic Oscillator Path Integral with MCMC}

\author{Javad Yousefian \and Alexandra VonHaven \and Kevin Geumhan}

\date{Version 1.3, November 14, 2026}

\begin{document}
\maketitle

\section{Introduction}
In this project, we plan to sample the path integral of the quantum harmonic oscillator in imaginary time using Markov chain Monte Carlo methods. We will write our own Metropolis sampler and also use \texttt{emcee}~\cite{emcee} so we can directly compare the two methods. This document describes our code validation plan. Version 1 was written before any code, so it does not include results; the later revisions are listed below.

Our main goal is to check if our Metropolis code samples the correct distribution. The harmonic oscillator is a good choice for validation because its lattice distribution is Gaussian, which lets us compare with exact analytical results.

\paragraph{Revisions.} Version 1 (November 9) was written before any code and without AI. Every later change is a separate commit in our Git repository, so the old text stays in the history and each commit message says what changed and why.
\begin{itemize}
  \item Version 1.1: \fillin{what changed and why}.
  \item Version 1.2 (November 14): our code samples a periodic lattice, not the fixed-endpoint lattice of version 1, so we rewrote Tests 6 and 8 to match the code, and made Test 11 precise.
  \item Version 1.3 (November 14): we rewrote Sections 3, 4, 5 and 7 for the periodic code, and Test 10, which now gets back both $m$ and $\omega$ and has a negative control. The parts that Kevin and Alexandra still update are marked in red.
\end{itemize}

\section{The model}
\fillin{Alexandra: rewrite this section for the periodic lattice our code uses (see Section 4 and Test 10).}

The Euclidean action of the oscillator is
\begin{equation}
  S_E = \int_0^{T_\tau} d\tau \left[ \frac{m}{2}\dot x^2 + \frac{m\omega^2}{2}x^2 \right].
\end{equation}
We divide the time into $N$ steps of size $\epsilon$ and use this discretization and rescaling:
\begin{equation}
  a = \omega\epsilon, \qquad \kappa = \frac{a^2}{4}, \qquad g = \frac{1-\kappa}{1+\kappa}, \qquad x_j = \sqrt{\frac{\hbar\epsilon}{m(1+\kappa)}}\, u_j .
\end{equation}
The endpoints are fixed, $u_0 = u_N = 0$, so the free variables are $u_1, \dots, u_{N-1}$. This is $D = N-1$ dimensions ($D = 63$ for our main lattice). The action becomes
\begin{equation}
  \mathcal S(u) = \sum_{j=1}^{N-1} u_j^2 - g \sum_{j=0}^{N-1} u_j u_{j+1},
  \label{eq:action}
\end{equation}
and we sample $\pi(u) \propto e^{-\mathcal S(u)}$. We keep the endpoints fixed and do not use periodic boundaries, because that would change the distribution.

\paragraph{Exact answer.} The action can be written as $\mathcal S = \frac12 u^T Q u$, where $Q$ has 2 on the diagonal and $-g$ next to the diagonal. Since $|g| < 1$, $Q$ is positive definite, so $\pi$ is a Gaussian with mean zero and covariance $Q^{-1}$:
\begin{equation}
  \langle u_j \rangle = 0, \qquad \langle u_j u_k \rangle = (Q^{-1})_{jk}, \qquad \langle u_j^4 \rangle = 3 \left[(Q^{-1})_{jj}\right]^2 .
  \label{eq:exact}
\end{equation}
We can get $Q^{-1}$ from NumPy, but there is also a closed formula for it. With $\mu = 2\,\mathrm{artanh}(a/2)$,
\begin{equation}
  (Q^{-1})_{jk} = \frac{\sinh(\mu \min(j,k))\, \sinh(\mu(N-\max(j,k)))}{g \sinh\mu \, \sinh(\mu N)} .
  \label{eq:Qinv}
\end{equation}
Far from the ends this is about $\frac{1+\kappa}{2a}\rho^{|j-k|}$ with $\rho = (2-a)/(2+a)$. For $a = 0.25$ this gives $\langle u_j^2\rangle \approx 2.03$, $\langle u_j u_{j+1}\rangle \approx 1.58$ and $\langle u_j u_{j+2}\rangle \approx 1.23$. Near the ends the variance is smaller: for $a = 0.25$ it is 0.802 at $j = 1$, 1.288 at $j = 2$ and 1.759 at $j = 4$.

\paragraph{Symmetries.} The distribution does not change if we flip the sign of the path ($u \to -u$) or if we reverse it in time ($u_j \to u_{N-j}$). So $\langle u_j\rangle = 0$ and $\langle u_j u_k\rangle = \langle u_{N-j}u_{N-k}\rangle$. There are no degeneracies, because $Q$ is positive definite and there is only one peak, at $u = 0$.

\paragraph{Energy gap.} In the middle of the lattice the correlation falls like $\rho^k$, so the decay rate per unit time is
\begin{equation}
  \frac{\omega_{\text{lat}}}{\omega} = \frac{2}{a}\,\mathrm{artanh}\!\left(\frac a2\right).
  \label{eq:wlat}
\end{equation}
For small $a$ this goes to 1, which is the continuum gap $E_1 - E_0 = \hbar\omega$. At $a = 1$ it is 1.099.

\section{What the code will do}
\begin{enumerate}
  \item Make a periodic path $x_0, \dots, x_{N-1}$, with $x_N = x_0$, and compute its action.
  \item Sample the path with our Metropolis code, and separately with \texttt{emcee}.
  \item Measure $\langle x\rangle$, $\langle x^2\rangle$, $\langle x^4\rangle$, the correlator $C(j) = \langle x_t x_{t+j}\rangle$ and the energies. The periodic lattice looks the same at every time slice, so we average over all slices $t$, which gives more statistics than using only the middle site.
  \item Fit the measured correlator with its exact lattice form to get back $m$ and $\omega$, and from them the energy gap (Test~10).
  \item Compute error bars and convergence checks.
  \item Save the settings, the seed, the results and the plots, so every run can be repeated. Every command of our package has an \texttt{-{}-out} option that does this, and \texttt{scripts/run\_all.sh} reproduces all our results.
\end{enumerate}

\section{Samplers}
\paragraph{Our Metropolis sampler.} We write the lattice action as $S = A\sum_t x_t^2 - B\sum_t x_t x_{t+1}$, with $A = m/a + am\omega^2/2$ and $B = m/a$. In one sweep we go through all $N$ sites of the periodic path; there are no fixed endpoints. For each site we propose $x_t' = x_t + u$, with $u$ uniform in $[-\delta, \delta]$. The change of the action only depends on the two neighbours:
\begin{equation}
  \Delta S = u\,[A(u + 2x_t) - B(x_{t-1}+x_{t+1})].
  \label{eq:dS}
\end{equation}
We accept with probability $\min(1, e^{-\Delta S})$. If we reject, we keep the old value and still count it as the next sample. To make it faster, we update all even sites at the same time and then all odd sites. This is allowed because even sites are not neighbours of each other; it needs $N$ to be even.

The expected acceptance rate can also be calculated exactly. Given its neighbours, $x_t$ is Gaussian with variance $1/(2A)$. If we measure the step in units where this variance is $1/2$, that is $\Delta = \delta\sqrt{A}$, averaging over this Gaussian gives
\begin{equation}
  P_{\text{acc}}(\Delta) = \mathrm{erfc}\!\left(\frac\Delta 2\right) + \frac{2}{\sqrt\pi\,\Delta}\left(1 - e^{-\Delta^2/4}\right).
  \label{eq:acc}
\end{equation}
For example $P_{\text{acc}}(1) = 0.729$, $P_{\text{acc}}(2) = 0.514$ and $P_{\text{acc}}(3) = 0.370$. In these units it does not depend on $a$, so the acceptance should be the same for every lattice spacing. In the code we set $\delta = \Delta/\sqrt{A}$.

We will try $\Delta = 0.5$, 1, 1.5, 2, 2.5, 3, and 4 in pilot runs (Test~8). After that, we will choose the value that gives the most effective samples with an acceptance rate between about 35\% and 65\%. We expect this value to be around 2.5, and we use the value chosen in Test~8 for all main runs.

\paragraph{emcee.} \fillin{Kevin: update this paragraph to match \texttt{ensemble.py}.} For the second sampler, we use \texttt{emcee.EnsembleSampler} with the same log probability $-\mathcal S(u)$ and at least $2D$ walkers. Because emcee works differently from our sampler, we expect the averages to match but not the individual chains. emcee is slow in high dimensions, so we will mainly use it for $N = 64$. If it is too slow, we will switch to $N = 16$ and note this change in the report. We will run emcee four times independently so we can compute $\hat R$.

\paragraph{Settings.} Our main lattice is $N = 64$ and $a = 0.25$, so $T = Na = 16$, with $m = \omega = 1$. For the injection test (Test~10) we also use $N = 16$ with $a = 1$, $N = 32$ with $a = 0.5$, and two cases with $m \ne 1$ and $\omega \ne 1$. We start four chains, each with 8,000 sweeps of burn-in and 40,000 sweeps for sampling, with the step chosen in Test~8. If the convergence checks fail, we will run longer. Two chains start at $x = 0$, and the other two start from a random path five times wider than the real distribution.

\section{Error bars and when a test passes}
Our samples are correlated, so we cannot use the normal $\sigma/\sqrt{M}$ error. We compute the integrated autocorrelation time
\begin{equation}
  \tau_{\text{int}} = \frac12 + \sum_{\ell=1}^{\ell_{\max}} \rho(\ell),
\end{equation}
where $\rho(\ell)$ is the autocorrelation at lag $\ell$ and $\ell_{\max}$ is chosen with Sokal's window~\cite{sokal}. The effective sample size is then $M/(2\tau_{\text{int}})$ for $M$ samples. For the error bars we average over blocks of consecutive sweeps, so that the correlations are included:
\begin{itemize}
  \item Tests 6 and 11 use binning: the error of the mean of the block averages, at the largest block size that still leaves at least 32 blocks.
  \item Test 10 uses 100 blocks, each at least $20\tau_{\text{int}}$ long, and also takes the covariance between the measured values from them for its correlated fit.
  \item The physics checks take the larger of two estimates: binning, averaged over the plateau, and $\sqrt{2\tau_{\text{int}}\,\mathrm{var}/M}$. For \texttt{emcee} they also use the spread between walkers. \fillin{Kevin: check this sentence.}
\end{itemize}

For the statistical tests we compare with the exact value using
\begin{equation}
  z = \frac{|\text{measured} - \text{exact}|}{\text{error bar}},
\end{equation}
and a test passes if $z \le 3.5$. We do many comparisons, so one of them could fail just by chance. If that happens, we run the test again with a different seed before we decide that there is a bug, and we report both runs. We will not change these limits after we see the results.

\section{Tests}
For every test we write what we check, when it passes, and what it would mean if it fails.

\fillin{Tests 1--5, 7 and 9 still describe the fixed-endpoint lattice; Kevin and Alexandra will update them.}

\subsection*{Test 1: Action and the update}
We compare $\Delta\mathcal S$ from Eq.~\eqref{eq:dS} with the difference of the full action for 1000 random paths. They should agree to $10^{-10}$. We also check that $\frac12 u^T Q u$ gives the same action, that the acceptance probability is always between 0 and 1, and that the fast odd/even sweep gives exactly the same result as a simple loop over the sites when we give both the same random numbers. If this fails, there is probably a sign or index mistake, or a problem at the endpoints.

\subsection*{Test 2: Exact covariance}
We check that our $Q$ is symmetric and positive definite, that its eigenvalues are $2 - 2g\cos(r\pi/N)$ for $r = 1, \dots, N-1$, and that the closed formula Eq.~\eqref{eq:Qinv} agrees with NumPy's inverse to $10^{-10}$. For $a = 0.25$ the numbers should match the values in Section 2. All the other tests compare with this, so if it fails we fix it first.

\subsection*{Test 3: Error bar code}
Before we trust our error bars, we test the autocorrelation and batch-means code on AR(1) data, $x_t = \phi x_{t-1} + \sqrt{1-\phi^2}\,\epsilon_t$, where the exact answer is $\tau_{\text{int}} = (1+\phi)/(2(1-\phi))$. We use $\phi = 0.5$, 0.9 and 0.99. The estimate of $\tau_{\text{int}}$ should be within 10\%, the batch-means error within 20\%, and emcee's value should be about $2\tau_{\text{int}}$. We also check that $\hat R$ is below 1.01 for four normal chains and above 1.01 if we shift one of them. If this fails, none of our error bars can be trusted.

\subsection*{Test 4: Very small lattices}
For $N = 2$ there is only one variable, and it should be Gaussian with variance $1/2$. For $N = 3$ there are two variables, and $u_1+u_2$ and $u_1-u_2$ should be independent Gaussians with variances $2/(2-g)$ and $2/(2+g)$. We run 200,000 sweeps and do a Kolmogorov--Smirnov test against these Gaussians, which passes if $p > 0.001$. The variance, the fourth moment and the acceptance rate should also pass $z \le 3.5$. This test can find mistakes that Test 1 cannot, for example if we forget to save the old value when a step is rejected.

\subsection*{Test 5: Known answer at $N = 64$}
This is the main test. With $N = 64$ and $a = 0.25$ we measure $\langle u_{j_0}\rangle$, $\langle u_{j_0}u_{j_0+k}\rangle$ for $k = 0, 1, 2, 4, 8$, $\langle u_{j_0}^4\rangle$, and the variance at the sites near the ends ($j = 1, 2, 4$ and $N-4$, $N-2$, $N-1$). All of them should pass $z \le 3.5$ against Eq.~\eqref{eq:exact}. If something fails, the sampler is biased, there is a bug near the ends, or our error bars are wrong.

\subsection*{Test 6 (revised): Symmetries and starting points}

\textbf{What changed.} On the periodic lattice there are no end sites, and there
is one more symmetry than before: shifting the path in time,
$x_t\to x_{t+s}$. Reversing the path in time is now $x_t\to x_{N-t}$. We also
state the pass limit of every check.

\textbf{What we check.} We use the four main chains of Section~4 ($N=64$,
$a=0.25$, $m=\omega=1$, $\Delta=2.5$; two start at $x=0$ and two from a random
path five times wider than the target; 8{,}000 burn-in and 40{,}000 kept sweeps
each). Each symmetry gives an exact prediction:
\begin{itemize}
\item $x\to-x$: $\langle x_t\rangle=0$ at $t=N/4$, $N/2$ and $3N/4$.
\item Time reversal: $\langle x_{j_0}x_{j_0+4}\rangle=\langle x_{j_0}x_{j_0-4}\rangle$
  with $j_0=N/2$, so their difference is 0.
\item Time shift: $\langle x_t^2\rangle=C(0)=0.4961$ at every site; we check
  $t=0$, $N/4$, $N/2$ and $3N/4$.
\end{itemize}
Error bars come from binning each chain and combining the four chains. Finally,
$\langle x^2\rangle$ and $C(4)$ from the two chains that start at $x=0$ must
agree with those from the two chains that start from the wide path.

\textbf{Degeneracies.} The distribution has none: $K$ is positive definite, so
there is a single peak at $x=0$. The only degeneracy is in the fit of Test~10,
where $\omega$ enters only as $\omega^2$, so $\omega$ and $-\omega$ fit equally
well; we take $\omega>0$.

\textbf{When it passes.} Each symmetry check passes $z\le3.5$, and the two
kinds of starting point agree within 3 error bars.

\textbf{If it fails.} The burn-in is too short, the chains do not mix well, or a
bug breaks a symmetry, for example a wrong index in the neighbour sum, which
would break time reversal.

\subsection*{Test 7: Convergence}
We look at trace plots of $u_{j_0}$, $u_{j_0}^2$ and the path average $\bar u$. The path average changes the slowest, so it is the most important one. For all the numbers we report, we need $\hat R < 1.01$ (the rank-normalized split $\hat R$ from~\cite{vehtari}) and an effective sample size above 400 in total. If this fails, we run longer or change $\Delta$.

\subsection*{Test 8 (revised): Acceptance rate and choice of step size}

\textbf{What changed.} Our code proposes $x_t\to x_t+u$ with $u$ uniform in
$[-\delta,\delta]$, in the units of $x$. Given its neighbours, $x_t$ is
Gaussian with variance $1/(2A)$, where $A=m/a+am\omega^2/2$, while
Eq.~\eqref{eq:acc} is written for variance $1/2$. A step $\Delta$ in the units
of Eq.~\eqref{eq:acc} is therefore $\delta=\Delta/\sqrt{A}$ in the code. With
this conversion Eq.~\eqref{eq:acc} holds unchanged and still does not depend
on $a$.

\textbf{What we check.} For each pilot step $\Delta=0.5,1,1.5,2,2.5,3,4$ we run
at $a=0.25$ ($N=64$, 40{,}000 sweeps) and at $a=1$ ($N=16$, 5{,}000 sweeps), with
$\delta=\Delta/\sqrt{A}$, and compare the measured acceptance, with a binning
error bar, to Eq.~\eqref{eq:acc}: 0.860, 0.729, 0.612, 0.514, 0.434, 0.370 and
0.282.

\textbf{When it passes.} Every measured acceptance agrees with
Eq.~\eqref{eq:acc} within $z\le3.5$, at both lattice spacings.

\textbf{Choice of $\Delta$.} At $a=0.25$ we also measure $\tau_{\rm int}$ of the
path average and of $x^2$, and choose the $\Delta$ that gives the most
effective samples per sweep, $1/(2\tau_{\rm int})$ for the slower of the two,
among the steps with acceptance between 35\% and 65\%. We expect about 2.5.

\textbf{If it fails.} There is a mistake in the $u^2$ part of $\Delta S$, in
the proposal width, or in the accept/reject step.

\subsection*{Test 9: Our sampler against emcee}
We compare $\langle u_{j_0}\rangle$, $\langle u_{j_0}^2\rangle$ and the correlation at $k = 1, 4, 8$ from our sampler and from emcee at $N = 64$. The difference divided by the combined error bar should be at most 3. Both samplers should also agree with the exact answer, because two samplers can agree with each other and still both be wrong. If one of them does not converge, we do not use it to judge the other one.

\subsection*{Test 10 (revised): Injection test --- getting back $m$ and $\omega$}

\textbf{What changed.} Version 1 of this test used the fixed-endpoint lattice
and its energy-gap formula. Our code samples the periodic lattice with the
action $S=\sum_t\big[\tfrac{m}{2a}(x_{t+1}-x_t)^2+\tfrac{a m\omega^2}{2}x_t^2\big]$,
so we rewrote the test for it, and we now recover both $m$ and $\omega$ instead
of only the energy gap.

\textbf{What we check.} We put known values of $m$ and $\omega$ into the
simulation, sample with our Metropolis code, and get $m$ and $\omega$ back from
the samples alone. On the periodic lattice the correlator is known exactly for
$0\le j\le N$:
\begin{equation}
C(j)=\langle x_t x_{t+j}\rangle
=\frac{\cosh\big(\mu(N/2-j)\big)}{2(m/a)\,\sinh\mu\,\sinh(\mu N/2)},
\qquad \cosh\mu=1+\frac{a^2\omega^2}{2}.
\end{equation}
We measure $C(j)$, averaged over all $t$, and fit this formula with $m$ and
$\omega$ as the two free parameters: the decay rate gives $\omega$ and the
overall size gives $m$. The fit uses times between 0.5 and 3 decay times, with
the decay rate estimated from the data, so the analysis never uses the injected
values. It is a correlated fit, with the covariance estimated from 100 blocks of
the chain, each at least $20\tau_{\rm int}$ long. The formula comes from the
action, not from our code, so a bug in the code cannot cancel out.

\textbf{Cases.} $(N,a,m,\omega)=(16,1,1,1)$, $(32,0.5,1,1)$, $(64,0.25,1,1)$,
$(32,0.5,2,0.5)$ and $(32,0.5,0.5,2)$, all with $T=Na=16$ and $10^5$ sweeps
each. The last two have $m\ne1$ and $\omega\ne1$, because with $m=\omega=1$ some
mistakes, such as a missing factor of $m$ in the potential, would not show up.

\textbf{When it passes.} For every case, the recovered $m$ and $\omega$ agree
with the injected values within 3.5 error bars, and the fit is good ($p>0.001$).

\textbf{Can the test fail?} As a negative control, we also sample with
$\omega=1.1$ but compare the result with $\omega=1$. The test must flag this
(pull above 3.5); if it does not, the test is too weak to be useful.

\textbf{Lattice effect.} The energy gap on this lattice is
$E_1-E_0=\mu/a=(2/a)\,\mathrm{arsinh}(a\omega/2)$, which is below $\omega$: it
is $0.962\,\omega$ at $a\omega=1$, $0.990\,\omega$ at $a\omega=0.5$ and
$0.997\,\omega$ at $a\omega=0.25$. We report it to show the discretization
effect, but the pass test is on $m$ and $\omega$ themselves.

\textbf{If it fails.} A wrong $\omega$ means the sampler gets the long-distance
correlations wrong, or the action in the code is not the one above. A wrong $m$
with a correct $\omega$ points to a wrong overall scale, for example a factor of
2 in the action. A bad fit means the measured correlator does not have the
expected shape, or our error bars are wrong.

\subsection*{Test 11 (revised): Multi-walker test}

\textbf{What changed.} We made ``roughly the same amount of time'' precise: the
test now states the interval, the number of walkers, how each walker's error
bar is found, and the pass limits, and it also compares with the exact fraction.

\textbf{What we check.} We run 8 independent walkers on the main lattice
($N=64$, $a=0.25$, $m=\omega=1$, $\Delta=2.5$): 4 start at $x=0$ and 4 from a
random path five times wider than the target distribution. After 8{,}000
burn-in sweeps, we record for 40{,}000 sweeps whether each walker is inside the
interval $|q|<\sigma_q$, for two quantities: $x$ at the middle site, with
$\sigma_q^2=C(0)=0.4961$, and the path average $\bar x$, with
$\sigma_q^2=1/(Nam\omega^2)=1/16$, which moves the slowest. Both quantities are
Gaussian, so the exact fraction of time inside is
$\mathrm{erf}(1/\sqrt2)=0.6827$. The error bar on each walker's fraction comes
from binning, so it includes the autocorrelation.

\textbf{When it passes.} For both quantities, every walker's fraction agrees
with 0.6827 within $z\le3.5$, and the walkers agree with each other: with
$\bar f$ the weighted mean of the fractions $f_i$,
$\chi^2=\sum_i (f_i-\bar f)^2/\sigma_i^2$ gives $p>0.001$ for 7 degrees of
freedom.

\textbf{If it fails.} If one walker differs from the others, it is stuck or its
burn-in was too short for its starting point. If the walkers agree with each
other but not with 0.6827, the sampler converges to the wrong distribution.

\section{Test execution}
We did not turn the tests into \texttt{pytest} files. Our tests are statistical runs whose numbers and figures we want to keep for the report, so instead every test runs with one command of our installed package, and \texttt{-{}-out} saves its output (numbers, PASS or FAIL, figures):
\begin{itemize}
  \item \texttt{sho-mcmc-validate} runs Tests 6, 8 and 11, and the other tests once they are added;
  \item \texttt{sho-mcmc-inject} runs Test 10;
  \item \texttt{sho-mcmc -{}-emcee} runs the physics and symmetry checks for both samplers (Test 9);
  \item \texttt{bash scripts/run\_all.sh} runs all of them and saves everything in \texttt{results/}.
\end{itemize}
The tests do not run automatically on each pull request. The author runs the relevant tests before opening a pull request, and the reviewer runs them again before merging. We have to look at the trace plots in Test~7 ourselves.

\section{Workflow}
We will add this plan to our Git repository before we start coding, so the Git history shows that the plan came before the code. After that, all changes (including code, tests, and this plan) will go through pull requests, and the other person will review them in writing before merging.

If a test fails or is not useful, we will update the plan in a new commit and keep the old version in the history. The commit message and the report will explain what changed and why.


\begin{thebibliography}{9}
\bibitem{emcee} D. Foreman-Mackey, D. W. Hogg, D. Lang and J. Goodman, ``emcee: The MCMC Hammer,'' PASP \textbf{125}, 306 (2013).
\bibitem{vehtari} A. Vehtari, A. Gelman, D. Simpson, B. Carpenter and P.-C. B\"urkner, ``Rank-normalization, folding, and localization: an improved $\hat R$ for assessing convergence of MCMC,'' Bayesian Analysis \textbf{16}, 667 (2021).
\bibitem{sokal} A. D. Sokal, \emph{Monte Carlo Methods in Statistical Mechanics: Foundations and New Algorithms}, lecture notes (1996).
\end{thebibliography}

\end{document}